%=====================================================================
\chapter{Dynamic Programming I: The Elements}\label{ch:dpone}
%=====================================================================
\locator{4}{Part IV --- The Design Paradigms}

\begin{outcomes}
By the end of this chapter you will be able to:
\begin{tight}
  \item \textbf{State} the two properties a problem must have for dynamic
        programming to apply, and \textbf{test} a given problem for them.
  \item \textbf{Distinguish} overlapping subproblems from the independent
        subproblems of divide-and-conquer.
  \item \textbf{Convert} a naive recursion into a memoised and then a bottom-up
        solution, and \textbf{compare} the three \emph{(CLO-6)}.
  \item \textbf{Derive} the running time of a dynamic program from the number
        of subproblems and the work per subproblem.
  \item \textbf{Decide} between greedy and dynamic programming for a stated
        problem.
\end{tight}
\end{outcomes}

\begin{prereq}
\chref{recurrences} for recursion trees, \chref{greedy} for optimal
substructure and the comparison that closes this chapter.
\end{prereq}

\begin{keyterms}
dynamic programming $\bullet$ optimal substructure $\bullet$ overlapping
subproblems $\bullet$ memoisation $\bullet$ top-down $\bullet$ bottom-up
$\bullet$ subproblem graph $\bullet$ rod cutting $\bullet$ reconstruction
$\bullet$ table
\end{keyterms}

\section{The Two Conditions}

\begin{definitionbox}[When Dynamic Programming Applies]
\textbf{1. Optimal substructure.} An optimal solution to the problem contains
optimal solutions to subproblems.

\smallskip
\textbf{2. Overlapping subproblems.} A recursive solution revisits the
\emph{same} subproblems repeatedly --- the total number of \emph{distinct}
subproblems is small, while the number of recursive calls is large.

\smallskip
\term{Dynamic programming} solves each distinct subproblem once and stores the
answer.
\end{definitionbox}

\begin{conceptbox}[Condition 2 is what makes it a technique]
Merge sort has optimal substructure: sorting an array optimally requires
sorting its halves. Memoising merge sort would buy exactly nothing, because
the two halves are \emph{different} subproblems and no subproblem ever recurs.

\smallskip
That is the dividing line:

\begin{tight}
  \item \textbf{divide and conquer} --- subproblems are independent and
        distinct, so recursion is the whole answer;
  \item \textbf{dynamic programming} --- subproblems overlap, so recursion
        alone does exponential redundant work.
\end{tight}

\smallskip
And unlike the greedy-choice property, overlap is \emph{measurable}: instrument
the recursion, count how many times each subproblem is solved, and look. That
is what Section~15.3 does, and the numbers are startling.
\end{conceptbox}

\section{Rod Cutting}

\begin{definitionbox}[Rod Cutting]
A rod of length $n$ can be cut into integer pieces. A piece of length $i$ sells
for $p_{i}$. Cut the rod to maximise total revenue (cutting is free).
\end{definitionbox}

Optimal substructure is easy to see: after the first cut of length $i$, the
remaining rod of length $n-i$ must itself be cut optimally, or the whole
solution could be improved. So, writing $r_{n}$ for the best revenue,
\[ r_{n} = \max_{1 \le i \le n} \bigl(p_{i} + r_{n-i}\bigr), \qquad r_{0} = 0. \]

\dsfig{\aoaalg{\algname{Cut-Rod-Naive}$(p, n)$}{
\If{$n = 0$}
  \State \Return $0$
\EndIf
\State $q \gets -\infty$
\For{$i \gets 1$ \textbf{to} $n$}
  \State $q \gets \max\bigl(q,\; p_{i} +
         \algname{Cut-Rod-Naive}(p, n-i)\bigr)$
    \Comment{re-solves $n-i$ every time}
\EndFor
\State \Return $q$
}}{\algname{Cut-Rod-Naive}: a direct transcription of the recurrence. It is
correct, and it takes $\bigTh{2^{n}}$ time, because line~5 re-solves
subproblems that other branches have already solved.}{cutrodnaive}

\begin{theorem}[The naive recursion is exponential]\label{thm:cutrodnaive}
\algname{Cut-Rod-Naive} makes exactly $2^{n}$ calls on an input of length $n$.
\end{theorem}

\begin{proof}
Let $T(n)$ be the number of calls, counting the top-level one. $T(0) = 1$, and
for $n \ge 1$ the body makes one call on each of $n-1, n-2, \ldots, 0$:
\[ T(n) = 1 + \sum_{j=0}^{n-1} T(j). \]
We show $T(n) = 2^{n}$ by induction. The base case holds, $T(0) = 1 = 2^{0}$.
Assuming it for all $j < n$,
\[ T(n) = 1 + \sum_{j=0}^{n-1} 2^{j} = 1 + (2^{n} - 1) = 2^{n}. \qedhere \]
\end{proof}

\begin{worked}[Measured: How Often Is Each Subproblem Solved?]
\texttt{code/ch15\_dp1.cpp} instruments the recursion and counts visits to each
subproblem. At $n = 22$:

\rowcolors{2}{white}{lightbg}
\begin{center}
\begin{tabular}{@{}rr@{}}
\toprule
\rowcolor{primary!15}
subproblem $k$ & times solved \\
\midrule
22 & 1 \\
21 & 1 \\
20 & 2 \\
19 & 4 \\
18 & 8 \\
17 & 16 \\
\multicolumn{2}{c}{$\cdots$ doubling all the way down $\cdots$} \\
 3 & 262{,}144 \\
 2 & 524{,}288 \\
 1 & 1{,}048{,}576 \\
 0 & \textbf{2{,}097{,}152} \\
\midrule
\textbf{total calls} & \textbf{4{,}194{,}304} \\
\bottomrule
\end{tabular}
\end{center}

\smallskip
\textbf{The subproblem ``a rod of length 0'' was solved 2{,}097{,}152 times.}
There is only one of it, and its answer is 0.

\smallskip
\textbf{Total calls $4{,}194{,}304 = 2^{22}$ exactly}, confirming
Theorem~\ref{thm:cutrodnaive} to the unit.

\smallskip
\textbf{And there are only 23 distinct subproblems} --- the rod lengths $0$ to
$22$. Four million calls to solve twenty-three problems. That ratio is the
saving dynamic programming is about to collect, and it is why ``overlapping
subproblems'' is a measurable condition rather than an aesthetic judgement.
\end{worked}

\section{Two Ways to Collect the Saving}

\begin{conceptbox}[Memoisation and bottom-up]
\textbf{Top-down with memoisation.} Keep the recursion exactly as written, and
add a table. Before computing, check whether the answer is known; after
computing, record it. One line of difference.

\smallskip
\textbf{Bottom-up.} Sort the subproblems by size and solve them in that order,
so every subproblem's dependencies are already in the table when it is reached.
No recursion at all.

\smallskip
\textbf{They have the same asymptotic running time}, and in practice:

\begin{tight}
  \item memoisation is a smaller edit to working code and only solves the
        subproblems actually needed --- which matters when the table is sparse;
  \item bottom-up has a smaller constant (no call overhead, no cache checks),
        touches memory in order, and \textbf{cannot overflow the stack} ---
        which for large $n$ is the deciding argument, as \chref{quicksort}
        showed.
\end{tight}
\end{conceptbox}

\dsfig{\aoaalg{\algname{Cut-Rod-DP}$(p, n)$}{
\State let $r[0 \mathinner{\ldotp\ldotp} n]$ be a new array
\State $r[0] \gets 0$
\For{$j \gets 1$ \textbf{to} $n$}
  \State $q \gets -\infty$
  \For{$i \gets 1$ \textbf{to} $j$}
    \State $q \gets \max\bigl(q,\; p_{i} + r[j-i]\bigr)$
      \Comment{$r[j-i]$ is already known}
  \EndFor
  \State $r[j] \gets q$
\EndFor
\State \Return $r[n]$
}}{\algname{Cut-Rod-DP}, bottom-up. Two nested loops, so $\bigTh{n^{2}}$ ---
and the only change from the recurrence is that subproblems are visited in
increasing order of size, so each is solved once.}{cutroddp}

\begin{theorem}[The dynamic program]\label{thm:cutroddp}
\algname{Cut-Rod-DP} computes $r_{n}$ in $\bigTh{n^{2}}$ time and
$\bigTh{n}$ space.
\end{theorem}

\begin{proof}
The outer loop runs $n$ times and the inner loop at most $j \le n$ times, with
an $\bigO{1}$ body, so the time is $\bigO{n^{2}}$; it is $\bigOm{n^{2}}$
because the inner loop runs exactly $j$ times, giving
$\sum_{j=1}^{n} j = n(n+1)/2$ iterations. The array $r$ gives $\bigTh{n}$
space.

\smallskip
Correctness is an induction on $j$: assuming $r[0], \ldots, r[j-1]$ hold the
optimal revenues for those lengths, line~6 evaluates $\max_{i}(p_{i} +
r_{j-i})$, which by the recurrence is $r_{j}$.
\end{proof}

\begin{conceptbox}[The running time formula worth memorising]
\[ \text{time} \;=\; (\text{number of distinct subproblems})
   \;\times\; (\text{work per subproblem}). \]
Rod cutting: $n$ subproblems $\times$ $\bigO{n}$ choices each $= \bigTh{n^{2}}$.
This formula handles nearly every dynamic program in \chref{dptwo}, and it is
usually quicker than writing down a recurrence.
\end{conceptbox}

\section{Measured}

\begin{worked}[The Explosion and the Collapse]
Calls made by each version:

\rowcolors{2}{white}{lightbg}
\begin{center}
\begin{tabular}{@{}rrrr@{}}
\toprule
\rowcolor{primary!15}
$n$ & naive calls & memoised calls & ratio \\
\midrule
10 &        1{,}024 &  56 &       18.3 \\
14 &       16{,}384 & 106 &      154.6 \\
18 &      262{,}144 & 172 &    1{,}524.1 \\
22 &    4{,}194{,}304 & 254 &   16{,}513.0 \\
26 &   67{,}108{,}864 & 352 & \textbf{190{,}650.2} \\
\bottomrule
\end{tabular}
\end{center}

\smallskip
The naive column is $2^{n}$; the memoised column grows like $n^{2}$. Their
ratio therefore grows like $2^{n}/n^{2}$ --- and does, from 18 to 190{,}000
over sixteen units of $n$.
\end{worked}

\begin{worked}[In Seconds]
\rowcolors{2}{white}{lightbg}
\begin{center}
\begin{tabular}{@{}rrrrr@{}}
\toprule
\rowcolor{primary!15}
$n$ & naive (ms) & memoised (ms) & bottom-up (ms) & naive/bottom-up \\
\midrule
20 &      4.113 & 0.0009 & 0.0004 &     10{,}281 \\
24 &     72.090 & 0.0012 & 0.0005 &    144{,}180 \\
28 &  1{,}130.981 & 0.0015 & 0.0006 &  1{,}884{,}968 \\
32 & 16{,}640.171 & 0.0017 & 0.0006 & \textbf{27{,}733{,}618} \\
\bottomrule
\end{tabular}
\end{center}

\smallskip
\textbf{At $n = 32$: 16.6 seconds against 0.6 microseconds.} A factor of 27.7
million, for a change that adds one array and one \texttt{if}.

\smallskip
\textbf{Bottom-up beats memoised by about 2.5$\times$} at every size --- the
constant-factor advantage of no recursion and no cache checks, exactly as
Section~15.3 claimed.

\smallskip
\textbf{And the naive column's ratios are 17.5, 15.7, 14.7} per four units of
$n$, against the predicted $2^{4} = 16$. The exponential is not a figure of
speech.
\end{worked}

\begin{worked}[Where the Naive Version Cannot Go]
Bottom-up alone, at sizes the recursion could not reach in the lifetime of the
universe:

\rowcolors{2}{white}{lightbg}
\begin{center}
\begin{tabular}{@{}rr@{}}
\toprule
\rowcolor{primary!15}
$n$ & bottom-up (ms) \\
\midrule
 10{,}000 & 0.263 \\
 40{,}000 & 1.054 \\
160{,}000 & 4.802 \\
\bottomrule
\end{tabular}
\end{center}

\smallskip
The times are \emph{linear} in $n$, not quadratic, because the price table here
is capped at 30 entries --- so the inner loop runs at most 30 times rather than
$j$ times, and $\bigTh{n^{2}}$ becomes $\bigTh{30n}$. A reminder to check what
the parameters actually are before quoting a bound: the algorithm is
$\bigTh{n \cdot \min(n, |p|)}$, and which factor dominates depends on the
instance.

\smallskip
At $n = 160{,}000$ the naive version would need about $2^{160000}$ calls.
\end{worked}

\section{Greedy or Dynamic Programming?}

\begin{conceptbox}[The comparison, finally]
\rowcolors{2}{white}{lightbg}
\begin{longtable}{@{}L{3.4cm}L{5.0cm}L{5.6cm}@{}}
\toprule
\rowcolor{primary!15}
& \textbf{Greedy} & \textbf{Dynamic programming} \\
\midrule
\endhead
needs & greedy-choice property \emph{and} optimal substructure & optimal
  substructure \emph{and} overlapping subproblems \\
decides & commits to one choice, never reconsiders & considers all choices,
  keeps the best \\
cost & usually $\bigTh{n}$ or $\bigTh{n\lg n}$ & usually
  $\bigTh{n^{2}}$ or worse \\
when it fails & silently returns a suboptimal answer & does not \\
proof obligation & an exchange argument & the recurrence is correct \\
\bottomrule
\end{longtable}

\smallskip
\textbf{Both require optimal substructure}; that is the shared half. They
differ in the second condition, and the second conditions are not opposites ---
a problem can have both, or neither.

\smallskip
\textbf{The practical rule:} try greedy first, because it is faster and
simpler. Attempt the exchange argument. If it goes through, you are done; if it
does not, the place it stuck usually hands you the counter-example, and dynamic
programming is the fallback that always works. Coin change in \chref{greedy} is
exactly this story: greedy for $\{1,5,10,25\}$, dynamic programming for
$\{1,3,4\}$.
\end{conceptbox}

\begin{alertbox}[Reconstructing the solution, not just its value]
Every algorithm in this chapter returns a \emph{number} --- the best revenue
--- not the cuts that achieve it. That is almost never what is wanted.

\smallskip
The fix is to store, alongside each $r[j]$, the choice $i$ that achieved the
maximum. Walking back from $j = n$ then produces the cuts. The cost is
$\bigTh{n}$ extra space and no extra time.

\smallskip
This matters more than it sounds, and \chref{dptwo} returns to it: the
space-saving tricks that reduce a DP table from $\bigTh{mn}$ to
$\bigTh{\min(m,n)}$ work only if you do \emph{not} need to reconstruct the
solution, because reconstruction needs the whole table.
\end{alertbox}

\begin{pitfall}
\begin{tight}
  \item \textbf{Memoising a problem with no overlap.} Merge sort gains nothing;
        the table is pure overhead.
  \item \textbf{Assuming optimal substructure.} It must be argued. Longest
        \emph{simple} path in a graph does not have it, and the usual DP
        recurrence for it is simply wrong.
  \item \textbf{Forgetting the base case in the table.} $r[0] = 0$ is load
        bearing.
  \item \textbf{Recursing memoised at large $n$.} The first call descends $n$
        frames deep before any memo is written --- \chref{quicksort}'s stack
        overflow, in a new costume. Bottom-up has no stack.
  \item \textbf{Returning only the value.} Store the choices too, or the
        answer is unusable.
  \item \textbf{Quoting $\bigTh{n^{2}}$ without checking the parameters.}
        Measured, rod cutting with a 30-entry price table is linear.
  \item \textbf{Reaching for DP before trying greedy.} Greedy is faster when it
        works, and the attempt to prove it tells you whether it does.
\end{tight}
\end{pitfall}

\begin{lab}[Counting the Overlap]
Reproduces all five tables. Expect twenty-five minutes; the naive version at
$n = 32$ takes about seventeen seconds on its own.

\begin{lstlisting}[language=C++]
// ch15_dp1.cpp -- overlapping subproblems, counted.
// Build: cl /O2 /EHsc /std:c++17 ch15_dp1.cpp        (see Appendix A)

// --- 1. The naive recursion, instrumented --------------------------
// Rod cutting: a rod of length n, a price p[i] for a piece of length
// i. Cut it to maximise revenue. The obvious recursion tries every
// first cut and recurses on the remainder.
static long long rod_naive(const std::vector<long long>& p, int n) {
    ++calls;
    if (n < (int)visits.size()) ++visits[n];
    if (n == 0) return 0;
    long long best = 0;
    for (int i = 1; i <= n && i < (int)p.size(); ++i)
        best = std::max(best, p[i] + rod_naive(p, n - i));
    return best;
}

// --- 2. The same recursion, with an answer cache -------------------
// One line of difference. Each subproblem is now solved once and
// looked up thereafter, which is why the running time collapses from
// exponential to quadratic.
static long long rod_memo(const std::vector<long long>& p, int n,
                          std::vector<long long>& memo) {
    ++calls;
    if (memo[n] >= 0) return memo[n];
    long long best = 0;
    if (n > 0)
        for (int i = 1; i <= n && i < (int)p.size(); ++i)
            best = std::max(best, p[i] + rod_memo(p, n - i, memo));
    return memo[n] = best;
}

// --- 3. Bottom-up: the same table, filled in dependency order ------
// No recursion and no cache lookups, just two loops. Identical
// answers, identical asymptotics, smaller constant -- and it cannot
// overflow the stack, which for large n matters (Chapter 9).
static long long rod_bottom_up(const std::vector<long long>& p, int n) {
    std::vector<long long> best(n + 1, 0);
    for (int j = 1; j <= n; ++j) {
        long long b = 0;
        for (int i = 1; i <= j && i < (int)p.size(); ++i)
            b = std::max(b, p[i] + best[j - i]);
        best[j] = b;
    }
    return best[n];
}

// --- 4. How often is each subproblem solved? -----------------------
// This is the measurement that DEFINES "overlapping". If every
// subproblem were solved once, DP would buy nothing.

// --- 5. The explosion, and the collapse ----------------------------
// --- 6. And in seconds ---------------------------------------------
// --- 7. All three agree, which is the point ------------------------
// --- 8. Bottom-up scales where the others cannot -------------------
\end{lstlisting}

\textbf{What to notice.} Part 4's last line: subproblem 0 solved 2{,}097{,}152
times. Every one of those calls returns 0 after doing nothing. The whole of
dynamic programming is the observation that this is avoidable.

\smallskip
\textbf{Then try to break it.} Remove the memo check from
\texttt{rod\_memo} but keep the memo \emph{write}. Predict what happens to the
call count before running it. The answer is that nothing improves at all ---
writing to a cache you never read is not memoisation, and the program will
still be exponential. The lookup is the technique; the table is just where it
keeps things.
\end{lab}

\begin{checkpoint}
\begin{tightnum}
  \item State the two conditions for dynamic programming and give a problem
        satisfying the first but not the second.
  \item \algname{Cut-Rod-Naive} made $2^{22} = 4{,}194{,}304$ calls to solve 23
        distinct subproblems. Derive the $2^{n}$ from the recurrence
        $T(n) = 1 + \sum_{j<n} T(j)$.
  \item Give the running time of a dynamic program with $mn$ subproblems each
        requiring a maximum over $k$ choices, and apply the formula to rod
        cutting.
\end{tightnum}
\end{checkpoint}

\begin{chaptersummary}
\begin{tight}
  \item Dynamic programming needs optimal substructure \emph{and} overlapping
        subproblems. The second distinguishes it from divide-and-conquer, where
        subproblems are distinct and memoising buys nothing.
  \item Overlap is measurable. Measured for rod cutting at $n = 22$: subproblem
        0 was solved 2{,}097{,}152 times, total calls were exactly $2^{22}$,
        and there were 23 distinct subproblems.
  \item Memoisation keeps the recursion and adds a table; bottom-up sorts the
        subproblems by size and uses loops. Same asymptotics; bottom-up has the
        smaller constant and no stack.
  \item Time $=$ (number of distinct subproblems) $\times$ (work per
        subproblem). Rod cutting: $n \times \bigO{n} = \bigTh{n^{2}}$.
  \item Measured: at $n = 32$ the naive version took 16.6 s and bottom-up
        0.0006 ms --- a factor of 27.7 million. Bottom-up beat memoised by
        about 2.5$\times$ throughout.
  \item Greedy and dynamic programming both need optimal substructure. Try
        greedy first and attempt the exchange argument; dynamic programming is
        the fallback that always works.
  \item Store the choices, not only the values, or the answer cannot be
        reconstructed.
\end{tight}
\end{chaptersummary}

\begin{reviewq}
  \item \textbf{(MCQ)} Dynamic programming requires: (a)~optimal substructure
        only; (b)~overlapping subproblems only; (c)~both; (d)~the
        greedy-choice property.
  \item \textbf{(MCQ)} Memoising merge sort would:
        (a)~make it $\bigTh{n}$; (b)~make it $\bigTh{\lg n}$; (c)~buy nothing,
        because its subproblems do not overlap; (d)~make it exponential.
  \item \textbf{(Short)} Explain the difference between overlapping and
        independent subproblems, with an algorithm of each kind.
  \item \textbf{(Short)} Compare memoised top-down with bottom-up on three
        axes, and say when each is preferable \emph{(CLO-6)}.
  \item \textbf{(Short)} Why does the space-saving trick of keeping only two
        rows of a DP table prevent reconstructing the solution?
  \item \textbf{(Applied)} Write the recurrence for coin change with
        denominations $c_{1},\ldots,c_{k}$ and amount $W$, state the number of
        subproblems and the work per subproblem, and give the running time.
  \item \textbf{(Applied)} A colleague memoises a recursive function and reports
        no speed-up at all. Give three distinct explanations, and the
        measurement that would distinguish them.
\end{reviewq}
